\section{Introduction}

An \(\qc{n,k,d}\) quantum code uses \(n\) physical qubits to store
\(k\) logical qubits. Its distance \(d\) is the smallest number of
qubits on which a Pauli operator can act nontrivially on the encoded
information while commuting with every stabilizer. Distance measures
the weight of the smallest undetected logical error. A code of distance
\(d\) corrects arbitrary errors on up to \(\lfloor(d-1)/2\rfloor\)
qubits. Both distance five and distance six correct two-qubit errors;
distance six also detects errors of weight five.

For fixed numbers of physical and logical qubits, determining the largest
possible minimum distance is a basic problem in quantum coding theory.
Constructions give lower bounds, while weight-enumerator identities and
other general bounds give upper bounds. When these differ by one, the
remaining existence question may depend on structure that an ordinary
weight enumerator does not record.

Let \(d_{\max}(n,k)\) denote the largest minimum distance of a binary
stabilizer code encoding \(k\) logical qubits into \(n\) physical qubits.
We prove the following two nonexistence results.
\begin{theorem}\label{thm:main}
There is no binary stabilizer code with parameters
\(\qc{16,5,d\geq5}\) or \(\qc{17,3,d\geq6}\). Hence
\[
 d_{\max}(16,5)=4,\qquad d_{\max}(17,3)=5.
\]
The conclusions hold for both pure and degenerate stabilizer codes.
\end{theorem}
The matching constructions are listed in Grassl's tables \cite{grassl}.
Their binary generators are in \pathcode{verification/data/}, and their distances
are checked by complete centralizer enumeration. The \(\qc{16,5,4}\)
construction extends a \(\qc{14,5,4}\) code by two coordinates;
the \(\qc{17,3,5}\) construction similarly extends a
\(\qc{15,3,5}\) code. The two exclusions close the displayed distance
intervals \(4\)--\(5\) and \(5\)--\(6\), respectively.

The proof begins by forcing a stabilizer word whose support leaves few
enough complementary coordinates for a complete classification. Let
\(T\) be that support and \(R\) the entire stabilizer supported on
\(T\). To shorten the centralizer on \(T\), retain only its words that
are identity on every coordinate of \(T\), then delete those coordinates.
Call the resulting additive code \(C\).
Every centralizer word has a component on \(T\) and a component on its
complement. Once the first component is fixed modulo \(R\), the second
belongs to a single coset of \(C\). The possible codes can be
described by linear maps between quotient spaces.

\begin{figure}[htbp]
\centering
\begin{tikzpicture}[>=Stealth,
 proofbox/.style={draw=black!60,fill=black!4,rounded corners=1pt,
 text width=4.25cm,minimum height=1.05cm,align=center,font=\small},
 node distance=0.65cm and 0.7cm]
\node[proofbox] (assume) {Assume the code exists};
\node[proofbox,right=of assume] (force) {Force a stabilizer word};
\node[proofbox,right=of force] (split) {Split the qubits on its support};
\node[proofbox,below=of split] (classify) {Classify the shortened codes};
\node[proofbox,left=of classify] (test) {Test every extension map};
\node[proofbox,left=of test] (none) {No map, so no code};
\draw[->] (assume)--(force);
\draw[->] (force)--(split);
\draw[->] (split)--(classify);
\draw[->] (classify)--(test);
\draw[->] (test)--(none);
\end{tikzpicture}
\caption{The structural reduction and the finite extension test.}
\label{fig:proof-flow}
\end{figure}

A first reading can follow the definitions in Section~\ref{sec:extension}
and the fourteen-qubit example in Section~\ref{sec:fourteen}, where the
counting inequality finishes the proof. Sections~\ref{sec:sixteen} and
\ref{sec:seventeen} apply the same setup to the two main targets;
Section~\ref{sec:computation} explains why their finite lists are complete.

\subsection{The extension criterion}

A coset of a subspace consists of all vectors obtained by adding its
members to one fixed vector. The quotient is the vector space of these
cosets: it treats two vectors as equivalent when their difference lies
in that subspace. Here this identifies local Pauli words that differ by
a stabilizer supported on the selected coordinates.

Put \(E=R_T^\perp/R\). Write \(c_R(e)\) for the minimum weight of a
representative of \(e\in E\), and write \(\delta(u)\) for the minimum
weight of a coset \(u\) of \(C\). In suitable coordinates, an extension
is described by a binary matrix \(M\) satisfying
\begin{equation}\label{eq:intro-criterion}
 \Gamma=MJM^\tr,\qquad c_R(e)+\delta(Me)\geq d\quad(e\ne0).
\end{equation}
Here \(J\) is the symplectic matrix on \(E\), and \(\Gamma\) is the
Gram matrix of a basis of \(C^\perp\). The first equation is the
commutation condition. The second states that the least weight in each
corresponding pair of cosets is at least \(d\).

Theorem~\ref{thm:extension} proves that these conditions are necessary
and sufficient when \(d(C)\geq d\) and every stabilizer word of weight
below \(d\) lies in \(R\). This includes pure codes and codes whose
degeneracy is confined to the selected coordinates.
Theorem~\ref{thm:relative} gives a relative-distance version without
this restriction. In that version, the relevant coset minima also depend
on \(M\).

For example, suppose \(R=\langle ZZZZ\rangle\) on four coordinates.
Here \(I\) denotes the identity and \(Z\) a nonidentity Pauli letter.
The local class \(ZIII+R=\{ZIII,IZZZ\}\) has minimum weight one.
If its image under \(M\) had coset minimum weight three, the two
parts could be combined into a centralizer word of weight four outside
\(R\). In the setting of Theorem~\ref{thm:extension}, this rules out
distance five. The inequality in \eqref{eq:intro-criterion} makes this
test for every local class at once.

Every nonzero element of \(\ker M\) represents a logical operator
supported entirely on \(T\). The weights and symplectic products of these operators restrict the hull
\(C\cap C^\perp\) of the shortened code. A further counting condition
compares cosets of \(C\) of small weight with classes in \(E/\ker M\)
of sufficiently large weight. These restrictions can exclude a shortened
code before its extension matrices are enumerated.

\subsection{Applications: the two nonexistence proofs}

For sixteen qubits, exact weight-enumerator certificates and an integer
congruence show that a hypothetical code is pure and has minimum
stabilizer weight \(m\in\{6,7,8,9\}\). Shortening the centralizer on
the support of such a word gives one of four additive-code classifications,
containing 3,957 equivalence classes altogether. None admits a map
satisfying the distance inequalities in \eqref{eq:intro-criterion}.
The commutation equation is not needed in this case: even this relaxation
has no solution.

For seventeen qubits, purity alone leaves a less convenient collection
of shortening cases. We instead prove that every hypothetical code has
a stabilizer word of weight eight. A quadratic-form argument bounds the
binary dimension of the shortened nine-coordinate code by six. It follows
that there are just five possible local stabilizers \(R\): one of rank
one and four of rank two. The rank-one type has five admissible kernel cases,
determined by the kernel dimension, restricted symplectic rank, and
shortened-code hull.
These restrictions reduce the remaining problem to complete classifications
on nine coordinates and finite extension searches.

As a smaller application, Section~\ref{sec:fourteen} gives the
fourteen-qubit argument in the same notation. There a unique weight-four
stabilizer word forces \(\ker M=0\). Of the 37 possible shortened
codes, 25 have the wrong hull dimension and the remaining 12 violate the
coset-counting condition. The final contradiction requires no search over
extension matrices.

\subsection{Previous work}

The symplectic representation of stabilizer codes and the MacWilliams and
shadow identities are standard \cite{crss,rains}. Gundersen et al.
\cite{puncturing} describe puncturing directly in symplectic and stabilizer
matrix notation. Complementary symplectic projections are treated in
\cite{anticodes}; related subspace constructions appear in the use of
Goursat's lemma for subsystem codes \cite{goursat} and in
entanglement-assisted error correction \cite{catalytic}.
Our extension criterion uses this linear algebra while retaining the
minimum weights in the original coordinate blocks. Its role here is to
derive the kernel and coset-counting restrictions and to specify the
finite problems used in the two nonexistence proofs.

Classification by extension and the use of automorphism groups are also
established methods \cite{kaski}. We use the distance-five additive-code
classifications of Grassl et al. \cite{dodeca}, and give completeness
checks for the additional classifications needed here. Earlier quantum
nonexistence proofs combining structural reductions and finite searches
include the \(\qc{13,5,4}\) result of Bierbrauer et al. \cite{bierbrauer}.
Related geometric methods for purity and local structure are developed
in \cite{geometry}.

The semidefinite-programming bounds with rational certificates in \cite{sdp}, Table~1,
leave code-space dimension \(K=32\) possible at \((n,d)=(16,5)\),
and \(K=8\) possible at \((17,6)\): their displayed intervals are
\(32\leq K\leq42\) and \(4\leq K\leq8\), respectively.
Theorem~\ref{thm:main} excludes the corresponding stabilizer codes by
using their additional algebraic structure.

Singh \cite{singh} announced the \(\qc{14,3,5}\) nonexistence result,
including degenerate codes, on 10 August 2026. Section~\ref{sec:fourteen}
gives a hull-and-coset argument as a worked application of the extension
criterion.

Lai \cite{lai} develops refined weight-enumerator and structural tests.
His Table~II records the intermediate candidate counts; Section~VII-D
leaves 27 candidate distributions for \(\qc{16,5,5}\) and seven for
\(\qc{17,3,6}\), with nonexistence unestablished for either parameter
set. The present proofs exclude both parameter sets through complete
shortened-code classifications and extension tests.

The computational premises, verification coverage, and data limits are
described in Section~\ref{sec:computation} and
Appendix~\ref{app:evidence}.

\subsection{Organization}

Section~\ref{sec:extension} proves the extension criterion.
Section~\ref{sec:obstructions} derives the kernel, hull, and coset-counting
bounds, including a redundancy bound when the local stabilizer is a
full-support plane. Sections~\ref{sec:fourteen}, \ref{sec:sixteen}, and
\ref{sec:seventeen} give the applications in increasing computational complexity. Section~\ref{sec:computation} describes classification
completeness and the computational verification. The appendices specify
the exact certificates, the extension algorithms, and the available data.
